\section{Locally nilpotent ideals}
\label{section-locally-nilpotent}

\noindent
Here is the definition.

\begin{definition}
\label{definition-locally-nilpotent-ideal}
Let $R$ be a ring. Let $I \subset R$ be an ideal.
We say $I$ is {\it locally nilpotent} if for every
$x \in I$ there exists an $n \in \mathbf{N}$ such
that $x^n = 0$. We say $I$ is {\it nilpotent} if
there exists an $n \in \mathbf{N}$ such that $I^n = 0$.
\end{definition}

\begin{example}
\label{example-locally-nilpotent-not-nilpotent}
Let $R = k[x_n | n \in \mathbf{N}]$ be the polynomial ring in infinitely
many variables over a field $k$. Let $I$ be the ideal generated by
the elements $x_n^n$ for $n \in \mathbf{N}$ and $S = R/I$. Then the ideal
$J \subset S$ generated by the images of $x_n$, $n \in \mathbf{N}$
is locally nilpotent, but not nilpotent. Indeed, since $S$-linear
combinations of nilpotents are nilpotent, to prove that $J$ is locally
nilpotent it is enough to observe that all its generators are nilpotent
(which they obviously are). On the other hand, for each $n \in \mathbf{N}$
it holds that $x_{n + 1}^n \not \in I$, so that $J^n \not = 0$.
It follows that $J$ is not nilpotent.
\end{example}

\begin{lemma}
\label{lemma-locally-nilpotent}
Let $R \to R'$ be a ring map and let $I \subset R$ be a locally nilpotent
ideal. Then $IR'$ is a locally nilpotent ideal of $R'$.
\end{lemma}

\begin{proof}
This follows from the fact that if $x, y \in R'$ are nilpotent, then
$x + y$ is nilpotent too. Namely, if $x^n = 0$ and $y^m = 0$, then
$(x + y)^{n + m - 1} = 0$.
\end{proof}

\begin{lemma}
\label{lemma-locally-nilpotent-unit}
Let $R$ be a ring and let $I \subset R$ be a locally nilpotent
ideal.
An element $x$ of $R$ is a unit if and only if the image of $x$
in $R/I$ is a unit.
\end{lemma}

\begin{proof}
If $x$ is a unit in $R$, then its image is clearly a unit in $R/I$.
It remains to prove the converse.
Assume the image of $y \in R$ in $R/I$ is the inverse of the image of $x$.
Then $xy = 1 - z$ for some $z \in I$.
This means that $1\equiv z$ modulo $xR$.
Since $z$ lies in the locally nilpotent ideal
$I$, we have $z^N = 0$ for some sufficiently large $N$.
It follows that $1 = 1^N \equiv z^N = 0$ modulo $xR$.
In other words, $x$ divides $1$ and is hence a unit.
\end{proof}

\begin{lemma}
\label{lemma-Noetherian-power}
\begin{slogan}
An ideal in a Noetherian ring is nilpotent if each element
of the ideal is nilpotent.
\end{slogan}
Let $R$ be a Noetherian ring. Let $I, J$ be ideals of $R$.
Suppose $J \subset \sqrt{I}$. Then $J^n \subset I$ for some $n$.
In particular, in a Noetherian ring the notions of
``locally nilpotent ideal''
and ``nilpotent ideal'' coincide.
\end{lemma}

\begin{proof}
Say $J = (f_1, \ldots, f_s)$.
By assumption $f_i^{d_i} \in I$.
Take $n = d_1 + d_2 + \ldots + d_s + 1$.
\end{proof}

\begin{lemma}
\label{lemma-lift-idempotents}
Let $R$ be a ring. Let $I \subset R$ be a locally nilpotent ideal.
Then $R \to R/I$ induces a bijection on idempotents.
\end{lemma}

\begin{proof}[First proof of Lemma \ref{lemma-lift-idempotents}]
As $I$ is locally nilpotent it is contained in every prime ideal.
Hence $\Spec(R/I) = V(I) = \Spec(R)$. Hence the
lemma follows from Lemma \ref{lemma-disjoint-decomposition}.
\end{proof}

\begin{proof}[Second proof of Lemma \ref{lemma-lift-idempotents}]
Suppose $\overline{e} \in R/I$ is an idempotent.
We have to lift $\overline{e}$ to an idempotent of $R$.

\medskip\noindent
First, choose any lift $f \in R$ of $\overline{e}$, and set
$x = f^2 - f$. Then, $x \in I$, so $x$ is nilpotent (since $I$
is locally nilpotent). Let now $J$ be the ideal of $R$ generated
by $x$. Then, $J$ is nilpotent (not just locally nilpotent),
since it is generated by the nilpotent $x$.

\medskip\noindent
Now, assume that we have found a lift $e \in R$ of $\overline{e}$
such that $e^2 - e \in J^k$ for some $k \geq 1$.
Let $e' = e - (2e - 1)(e^2 - e) = 3e^2 - 2e^3$, which is another
lift of $\overline{e}$ (since the idempotency of $\overline{e}$
yields $e^2 - e \in I$). Then
$$
(e')^2 - e' = (4e^2 - 4e - 3)(e^2 - e)^2 \in J^{2k}
$$
by a simple computation.

\medskip\noindent
We thus have started with a lift $e$ of $\overline{e}$ such
that $e^2 - e \in J^k$, and obtained a lift $e'$ of
$\overline{e}$ such that $(e')^2 - e' \in J^{2k}$.
This way we can successively improve the approximation
(starting with $e = f$, which fits the bill for $k = 1$).
Eventually, we reach a stage where $J^k = 0$, and at that
stage we have a lift $e$ of $\overline{e}$ such that
$e^2 - e \in J^k = 0$, that is, this $e$ is idempotent.

\medskip\noindent
We thus have seen that if $\overline{e} \in R/I$ is any
idempotent, then there exists a lift of $\overline{e}$
which is an idempotent of $R$.
It remains to prove that this lift is unique. Indeed, let
$e_1$ and $e_2$ be two such lifts. We
need to show that $e_1 = e_2$.

\medskip\noindent
By definition of $e_1$ and $e_2$, we have $e_1 \equiv e_2
\mod I$, and both $e_1$ and $e_2$ are idempotent. From
$e_1 \equiv e_2 \mod I$, we see that $e_1 - e_2 \in I$,
so that $e_1 - e_2$ is nilpotent (since $I$ is locally nilpotent).
A straightforward
computation (using the idempotency of $e_1$ and $e_2$)
reveals that $(e_1 - e_2)^3 = e_1 - e_2$. Using this and
induction, we obtain $(e_1 - e_2)^k = e_1 - e_2$ for any
positive odd integer $k$. Since all high enough $k$ satisfy
$(e_1 - e_2)^k = 0$ (since $e_1 - e_2$ is nilpotent),
this shows $e_1 - e_2 = 0$, so that $e_1 = e_2$, which
completes our proof.
\end{proof}

\begin{lemma}
\label{lemma-lift-idempotents-noncommutative}
Let $A$ be a possibly noncommutative algebra.
Let $e \in A$ be an element such that $x = e^2 - e$ is nilpotent.
Then there exists an idempotent of the form
$e' = e + x(\sum a_{i, j}e^ix^j) \in A$
with $a_{i, j} \in \mathbf{Z}$.
\end{lemma}

\begin{proof}
Consider the ring $R_n = \mathbf{Z}[e]/((e^2 - e)^n)$. It is clear that
if we can prove the result for each $R_n$ then the lemma follows.
In $R_n$ consider the ideal $I = (e^2 - e)$ and apply
Lemma \ref{lemma-lift-idempotents}.
\end{proof}

\begin{lemma}
\label{lemma-lift-nth-roots}
Let $R$ be a ring. Let $I \subset R$ be a locally nilpotent ideal.
Let $n \geq 1$ be an integer which is invertible in $R/I$. Then
\begin{enumerate}
\item the $n$th power map $1 + I \to 1 + I$, $1 + x \mapsto (1 + x)^n$
is a bijection,
\item a unit of $R$ is a $n$th power if and only if its image in $R/I$
is an $n$th power.
\end{enumerate}
\end{lemma}

\begin{proof}
Let $a \in R$ be a unit whose image in $R/I$ is the same as the image
of $b^n$ with $b \in R$. Then $b$ is a unit
(Lemma \ref{lemma-locally-nilpotent-unit}) and
$ab^{-n} = 1 + x$ for some $x \in I$. Hence $ab^{-n} = c^n$ by
part (1). Thus (2) follows from (1).

\medskip\noindent
Proof of (1). This is true because there is an inverse to the
map $1 + x \mapsto (1 + x)^n$. Namely, we can consider the map
which sends $1 + x$ to
\begin{align*}
(1 + x)^{1/n}
& =
1 + {1/n \choose 1}x +
{1/n \choose 2}x^2 +
{1/n \choose 3}x^3 + \ldots \\
& =
1 + \frac{1}{n} x + \frac{1 - n}{2n^2}x^2 +
\frac{(1 - n)(1 - 2n)}{6n^3}x^3 + \ldots
\end{align*}
as in elementary calculus. This makes sense because the series is finite
as $x^k = 0$ for all $k \gg 0$ and each coefficient
${1/n \choose k} \in \mathbf{Z}[1/n]$ (details omitted; observe that
$n$ is invertible in $R$ by Lemma \ref{lemma-locally-nilpotent-unit}).
\end{proof}






